% !PRINTED 538
\weberpage{538}
%<<CONT
rangement and their initial element, it follows that by the application of any permutation $\pi_1$ from $P$ the elements of the individual cycles $\gamma$ are not separated from one another, but can only be interchanged among themselves (cyclically), and moreover the cycles can be interchanged with one another. The group is therefore, if $s>1$, \gesperrt{imprimitive}.

A rational function of the arguments $\alpha,\alpha_1,\ldots\alpha_{r-1}$ which does not change its value when the $\alpha$ are subjected to the cyclic permutation $\gamma$ and its repetitions is called a \gesperrt{cyclic function} of the $\alpha$. If one denotes by
\weq{4}{\omega = \psi(\alpha,\alpha_1\ldots\alpha_{r-1})}
such a cyclic function, and by
\weq{5}{\omega,\ \omega_1,\ \omega_2\ldots\omega_{s-1}}
the conjugate values of $\omega$, setting thus e.\,g.
\weq{6}{\omega_1 = \psi(\beta,\beta_1\ldots\beta_{r-1}),}
then these quantities are, by \S.\,158, the roots of an irreducible equation of degree $s$:
\weq{7}{\Phi(t) = 0,}
whose group one obtains by seeking out the permutations of the quantities (5) produced by $P$.

But one obtains the permutation effected by $\pi_1\pi_2$ among the quantities (5) by composing the two permutations produced by $\pi_1$ and $\pi_2$ singly, and the group of the permutations of (5) is therefore also commutative. Hence $\Phi(t)=0$ is an \gesperrt{Abelian} equation of degree $s$.

If one adjoins $\omega$, then $F(x)$ splits into $s$ factors of degree $r$:
\[
F(x) = F(x,\omega)F(x,\omega_1)\ldots F(x,\omega_{s-1}),
\]
of which the first, $F(x,\omega)$, has the roots $\alpha,\alpha_1\ldots\alpha_{r-1}$, and the group of the equation $F(x,\omega)=0$ consists solely of the period of the cyclic permutation $\gamma$.

We shall call an equation whose group consists of a single cycle and its repetitions a \gesperrt{cyclic equation}, so that the cyclic equations are the simplest special case of the \gesperrt{Abelian} equations. We have then proved that the solution of every \gesperrt{Abelian} equation is reduced to the solution of an \gesperrt{Abelian} equation of lower degree and to the solution of a series
%CONT>>
